% Template for post or presentation abstracts submitted to ILAS 2022
% 1. Fill in the presenter's information (name, affiliation, email
% address) and talk information (title of presentation, abstract).
% 2. Compile to make sure there are no errors.
% 3. Save the .tex file for your abstract with a distinctive name,
% ideally "FamilynameGivenname.tex".
% 4. Upload your .tex file to https://clr.ie/129557
%       
% * Please keep your LaTeX commands simple, and avoid the use of
%    user-defined macros.
% * Include details of co-authors and funding acknowledgements after the abstract.
% * Upload only the LaTeX file (with .tex extension).
% * Questions: Contact Niall Madden (Niall.Madden@NUIGalway.ie)

\documentclass[a4paper]{amsart} 
\usepackage{amssymb} 
\usepackage{hyperref}
\pagestyle{empty}
\date{} 

%% IGNORE THE NEXT 4 LINES
\makeatletter % 
\def\labelenumi{\theenumi.}
\def\theenumi{\@arabic\c@enumi}
\makeatother
%% STOP IGNORING NOW

\begin{document}

\title{Rational matrix solutions to $p(X)=A$}
\author{Andr\'e Ran} %Presenting author only
\address{Vrije Universiteit, The Netherlands and North-West University, South Africa} % Affiliation of presenting author
\email{a.c.m.ran@vu.nl} % Email address of presenting author only

\maketitle

% The abstract  ***no more than one page***
Consider an $n\times n$ matrix $A$ with rational entries, and let $p(\lambda)$ be a polynomial with rational coefficients. The question we consider is whether or not there is a rational $n\times n$ matrix $X$ such that $p(X)=A$, and if there is, how to find it. In addition, we consider the problem over which field extension of the rationals there will be a solution, provided a solution over the complex numbers does exist.

The solution to this problem was first discussed in \cite{Drazin_paper} for the case where $A$ has $n$ distinct eigenvalues. Our contribution is to extend this result to more general cases. This requires a new canonical form for the matrix $A$.

In the talk we will outline the steps involved in the solution of the problem, and if there is time, discuss some examples. In particular, we are interested in $m$th roots of $A$.


%% Coauthor details here (optional - delete if not applicable)
% and funding acknowledgment (optional - delete if not applicable)
\emph{This is joint work with Gilbert Groenewald, Dawie Janse van Rensburg, Madelein van Straaten and Frieda Theron (all North-West University)}

\begin{thebibliography}{1}
\bibitem{Drazin_paper} 
M.P.~Drazin. 
\newblock Exact rational solutions of the matrix equation $A=p(X)$ by linearization.
\newblock {\em Linear Algebra Appl.}, 426 (2007) 502--515. 
\end{thebibliography}


\end{document}


% Please leave the next bit alone  
%%% Local Variables: 
%%% mode: latex
%%% TeX-master: "../ILAS-2022"
%%% End:
